% External validation against published methods (replication programme, 2026-11).
% Numbers: docs/REPLICATION_RESULTS.md (final, 2026-11-21) and the files it cites under results/replication/;
% rules: docs/PREREGISTRATION_REPLICATION.md (A1-A43) and docs/PREREGISTRATION_PHASE5.md.
\section{External validation against published methods}\label{sec:replication}
The results so far compare masking with remedies we implemented. To test the location law and the class-conditional
constraint against the field, we reproduced eight published shortcut-mitigation papers in their own environment (their
code, data, splits and metric) and ran their best methods in ours (thyroid calipers, Trap~A and Trap~B), under rules
registered before each result.

\paragraph{Protocol} For each paper we first ran its best method with the authors' code and compared it with the
published number (the \emph{reproduction gate}; tolerance 0.05 in the paper's metric, widened from 0.03 for every paper
after the first gates were known, which we report as a deviation). Where the gate passed (\emph{tier~1}), a method of
ours counts as having beaten theirs only if its mean is higher and the 95\% CI of the paired difference, Holm-adjusted
within the paper's family, lies above~0. Where the gate failed (\emph{tier~2}), we compare against our run of the
authors' official code under identical conditions and make no claim against the published number. Four variants of our
method were registered in advance: (a) the constraint on features of the competitor's own trained network; (b) the
constraint on frozen DINOv2 features of ROI-masked images, our pipeline; (c) the label-free U-MtE; (d) the constraint on
an EfficientNet-B0 trained on ROI-masked images (5.3\,M parameters, smaller than every competitor's network). The
competitors' artifact labels, ROIs and selection rules were used unchanged; heads were selected on validation data only,
and every test set was scored once.

\begin{table}[t]\centering\small
\caption{Replication programme: summary. Tier 1 = the paper's best method reproduced (verdicts apply); tier 2 = not
reproduced (descriptive comparison with our run of the authors' code). AUROC unless stated; $\Delta$ = our best
registered variant minus their method; 95\% CIs Holm-adjusted where a verdict applies.}\label{tab:repl_summary}
\begin{tabular}{>{\raggedright\arraybackslash}p{3.6cm}>{\raggedright\arraybackslash}p{2.5cm}>{\raggedright\arraybackslash}p{1.3cm}>{\raggedright\arraybackslash}p{4.4cm}>{\raggedright\arraybackslash}p{2.8cm}}\toprule
Paper (artifact) & Reproduced? & Tier & Ours vs theirs & Reading\\\midrule
Phase 5: ISNet, RR-ClArC on thyroid calipers inside the ROI & --- (our design) & confirm. & (b) $+0.155$ vs ISNet, $+0.249$ vs RR-ClArC & beaten, 4/4 \\
EPVT, MICCAI 2023 / TMI 2025 (hair, rulers, gel; trap sets) & yes (gap 0.050 / 0.017) & 1 & (b) $+0.086$ [$+0.064,+0.107$] (bias 1), $+0.031$ [$+0.009,+0.054$] (bias 0.9) & beaten at both levels\\
DISC, ICML 2023 (hair; trap sets) & no (0.489 vs 0.551) & 2 & (b) 0.736 vs 0.489, $+0.247$ [$+0.204,+0.290$] & above official code and published value\\
RR-ClArC, Mach.\ Learn.\ 2025 (inserted vignette) & no (0.524 vs 0.76, accuracy) & 2 & (d) 0.785 vs 0.524 & above official code and published value\\
Viviano et al., ICLR 2021: location test & --- (our test) & formal, 3 seeds & (a), (b) gain inside $+0.20$, CIs reach $-0.006$ & masking's crossover supported; gain not supported\\
ISNet, Nat.\ Commun.\ 2024 (marks) & no (external test) & 2 & see Table~\ref{tab:repl_isnet} & ISNet fails inside the ROI\\
Saab et al., MLHC 2022 (trap sets) & no (0.647 vs 0.727) & 2 & (b) $+0.023$ [$-0.029,+0.076$] & tie\\
RR-ClArC, AAAI 2024 (invisible whole-image trigger) & yes & 1 & (a) $-0.004$ (ResNet-18) & not beaten (out of scope)\\
\bottomrule\end{tabular}\end{table}

\paragraph{Confirmatory test on real calipers} In our environment we trained ISNet \citep{bassi2024isnet} and RR-ClArC
\citep{pahde2025ensuring} with their published recipes on the thyroid traps (3 folds; ISNet from ImageNet weights,
RR-ClArC on a ResNet-50d base with $\lambda$ and layer chosen on clean validation data), beside the constraint on frozen
DINOv2 features and on a trained EfficientNet-B0 (Table~\ref{tab:repl_phase5}). With the caliper inside the nodule, both
variants of the constraint beat both published methods in all four registered comparisons (Holm over four, crossed fold
$\times$ image bootstrap), and they were also better on clean images. ISNet, which penalises relevance outside the ROI,
was better when the caliper lay outside the nodule (0.673) than inside it (0.561), the pattern the location law predicts
for any method that acts only outside the ROI; its Trap~A result varied strongly across folds (0.681, 0.637, 0.365).

\begin{table}[t]\centering\small
\caption{Confirmatory comparison on thyroid calipers (TN3K; registered before any result). Reversed-test AUROC, mean
$\pm$ SD over three folds; $\Delta$ vs each competitor with 95\% CI (Holm over four).}\label{tab:repl_phase5}
\begin{tabular}{lccc}\toprule
Method & Trap A (inside ROI) & Trap B (outside) & Clean (Trap A)\\\midrule
Constraint, frozen DINOv2 (ours) & \textbf{0.716} $\pm$ 0.027 & 0.789 & 0.797\\
Constraint, EfficientNet-B0 (ours) & 0.690 $\pm$ 0.076 & 0.746 & 0.733\\
ISNet \citep{bassi2024isnet} & 0.561 $\pm$ 0.172 & 0.673 & 0.633\\
RR-ClArC \citep{pahde2025ensuring} & 0.467 $\pm$ 0.049 & 0.595 & 0.643\\
U-MtE (ours, label-free) & 0.504 & 0.782 & 0.727\\
ERM (ResNet-50d) & 0.440 & 0.546 & 0.625\\\midrule
\multicolumn{4}{l}{DINOv2 vs ISNet $+0.155$ [$+0.019,+0.316$]; vs RR-ClArC $+0.249$ [$+0.192,+0.304$]}\\
\multicolumn{4}{l}{EfficientNet-B0 vs ISNet $+0.129$ [$+0.035,+0.239$]; vs RR-ClArC $+0.224$ [$+0.159,+0.292$]}\\
\bottomrule\end{tabular}\end{table}

\paragraph{Their environment: artifacts on the lesion} EPVT \citep{yan2023epvt,yan2025pldg} treats dermoscopy artifact
types as domains and is evaluated on Bissoto's trap sets \citep{bissoto2022artifact}, in which hair, rulers and gel
often lie on the lesion. After a full audit of its code against the paper (37 items; one port error corrected, and the
paper's mixup followed where the released code differs), EPVT reproduced within the widened tolerance (five seeds;
bias~1: 0.696 vs 0.746 read from the published figure, gap 0.0497, smaller than the reading uncertainty; bias~0.9:
0.775 vs 0.792). The constraint on frozen DINOv2 features beat EPVT at both bias levels (0.782 and 0.806; Holm over
eight, crossed seed $\times$ image bootstrap), and the constraint applied to EPVT's own features also improved on EPVT
(0.758 and 0.795). EPVT was not compared with a clean condition in the registration; on the artifact-free bias-0 test
the frozen-encoder variant was 0.02--0.04 AUROC below EPVT, whereas the trained EfficientNet-B0 variant was within 0.01
and beat EPVT at bias~1 only.

DISC \citep{wu2023disc} did not reproduce on the same benchmark: with one image cache per trap set, ERM and DISC scored
0.507 and 0.489 against the published 0.364 and 0.551. The released code writes one image cache shared by all trap sets,
so that every trap set after the first is read in the first set's images; our misaligned ERM runs (0.368) matched the
published ERM value, which suggests that the published numbers carry this issue (we have informed the authors). Under
identical conditions all four variants of ours exceeded DISC run from its official code (DINOv2: 0.736 vs 0.489, $+0.247$
[$+0.204, +0.290$], higher on each of the five trap sets, and also when selection used no test-distribution images), and
the DINOv2 variant also exceeded the published DISC (0.551) and CBDebug \citep{enouen2026cbdebug} (0.580) values on this
benchmark. On the ISIC 2019 vignette of \citet{pahde2025ensuring}, RR-ClArC did not reproduce (0.524 vs 0.76 biased-test
accuracy); the EfficientNet-B0 variant reached 0.785 [0.769, 0.801], above both, at 0.807 clean accuracy (0.86 published
for RR-ClArC).

\paragraph{Published masking methods fail inside the ROI} In ISNet's own pipeline \citep{bassi2024isnet} we moved its
class marks from the image corner into the lungs (Table~\ref{tab:repl_isnet}). With the marks in the corner, ISNet ignored
them (macro-F1 0.84 / 0.83 / 0.83 on in-distribution images with own-class, no and swapped marks); with the same marks
inside the lungs, it relied on them alone (0.998 / 0.167 / 0.000). The published external-test level did not reproduce
for any model, including the authors' released weights, so these comparisons are descriptive. Our variants were
also driven largely by the marks when every class carried its own mark shape inside the ROI, although the frozen-encoder
variants retained more signal than ISNet under swapped marks (0.20--0.36 vs 0.03): when the artifact identifies the class,
removing within-class artifact directions cannot remove it. In the chest-radiograph setup of \citet{viviano2021saliency},
a radio-opaque tube moved inside the lungs+heart ROI lowered their masked model to 0.615 reversed-test AUROC (0.896 with
the tube outside; three seeds), whereas the constraint scored 0.867--0.876 at every seed. The registered gain over their
standard model was not significant (CIs reaching $-0.006$), because at one seed their checkpoint rule kept a model that
had not yet learned the tube.

\begin{table}[t]\centering\small
\caption{ISNet with its marks outside (corner) and inside the lungs (macro-F1; own-class / no / swapped marks).
Descriptive: in-distribution validation images (used for model selection) and, in brackets, the rebuilt external
test.}\label{tab:repl_isnet}
\begin{tabular}{llccc}\toprule
Marks & Method & Own-class & None & Swapped\\\midrule
Corner, all images & ISNet & 0.841 (0.427) & 0.834 (0.385) & 0.831 (0.387)\\
Corner, all images & U-MtE (ours) & 0.848 (0.383) & 0.848 (0.383) & 0.848 (0.383)\\
Inside lungs, all images & ISNet & 0.998 (0.999) & 0.167 (0.160) & 0.000 (0.000)\\
Inside lungs, 90/10 & ISNet & 0.974 (0.989) & 0.745 (0.435) & 0.033 (0.034)\\
Inside lungs, 90/10 & Constraint, DINOv2 & 0.980 (0.947) & 0.798 (0.398) & 0.197 (0.136)\\
Inside lungs, 90/10 & U-MtE & 0.973 (0.837) & 0.780 (0.389) & 0.359 (0.168)\\
Inside lungs, 90/10 & Constraint, EfficientNet-B0 & 0.982 (0.993) & 0.813 (0.406) & 0.033 (0.002)\\
\bottomrule\end{tabular}\end{table}

\paragraph{Where the constraint did not win} Segmentation-supervised ERM \citep{saab2022spatial}, which uses the lesion
masks as training targets, did not reproduce (0.647 vs 0.727) and tied with our best variant under identical conditions
($+0.023$ [$-0.029, +0.076$]); adding the constraint to its encoder lowered it ($-0.095$). RR-ClArC on an invisible
trigger written into the low bits of every pixel \citep{dreyer2024hope} reproduced and was not beaten; that artifact has
no location, so the setting lies outside the question of this paper. One further paper \citep{pahde2023r2r} did not
pass its reproduction check and was not compared. We emailed the authors of every paper that did not reproduce.
